\section{Introduction}

% L'introduction doit a minima presenter :
%   - le probleme etudie, plus en detail que dans le resume ;
%   - les idees principales de la modelisation ;
%   - le plan suivi dans le rapport.
% Un expose du contexte scientifique sur les approches possibles est un PLUS,
% avec les references bibliographiques correspondantes.

\subsection{Le probleme etudie}

% Evacuation d'urgence d'une piscine couverte. Decrire les trois milieux et
% pourquoi leur enchainement est interessant.

\subsection{Contexte scientifique}

% Les deux grandes familles de modeles de foule : approches continues (la foule
% comme un fluide) et approches particulaires (chaque individu est un objet).
% Au sein des approches particulaires : forces sociales (Helbing) contre
% contraintes dures (Maury et Venel). Dire pourquoi nous retenons la seconde.
% Introduire le diagramme fondamental comme observable de reference.

\subsection{Plan du rapport}

% A rediger en dernier.
